%
% Documento: Resumo (Inglês)
\begin{RESUMO}
\thispagestyle{empty}

	Este trabalho aborda a inclusão de crianças neurodivergentes no ensino fundamental público, destacando as dificuldades enfrentadas por professores devido à falta de formação continuada e recursos adaptados. O objetivo central é apresentar o IncluiEdu, um sistema web projetado para apoiar docentes na criação de atividades pedagógicas inclusivas. A metodologia baseou-se em uma abordagem de desenvolvimento de software estruturada em etapas de definição de requisitos, modelagem do sistema e implementação de uma arquitetura em nuvem, utilizando o Cloudflare Workers, o banco de dados relacional D1 e a Inteligência Artificial Generativa com o modelo Ollama. Como resultado, a plataforma permite o cadastro de alunos, o acompanhamento pedagógico contínuo e a geração automática de atividades adaptadas ao perfil do estudante. Para garantir a segurança e a ética, todo conteúdo gerado pela IA é salvo como rascunho, exigindo a revisão final do professor. Conclui-se que o IncluiEdu atua como uma ferramenta viável de apoio à decisão docente, reduzindo a sobrecarga no planejamento de aulas e contribuindo para uma educação mais equitativa, sem substituir a autonomia do educador.
	
	\vspace*{0.5cm}\noindent\textbf{Palavras-chave}: Deficiência intelectual; Educação inclusiva; Inteligência artificial.
	
\end{RESUMO}

\begin{ABSTRACT}
	
	This work addresses the inclusion of neurodivergent children in public elementary education, highlighting the difficulties faced by teachers due to the lack of continuous training and adapted resources. The main objective is to present IncluiEdu, a web system designed to support educators in creating inclusive pedagogical activities. The methodology was based on a software development approach structured in stages of requirements definition, system modeling, and implementation of a cloud architecture using Cloudflare Workers, the D1 relational database, and Generative Artificial Intelligence with the Ollama model. As a result, the platform allows student registration, continuous pedagogical monitoring, and the automatic generation of activities adapted to the student's profile. To ensure safety and ethics, all content generated by AI is saved as a draft, requiring the teacher's final review. It is concluded that IncluiEdu acts as a viable tool to support teaching decisions, reducing the workload in lesson planning and contributing to a more equitable education, without replacing the educator's autonomy.

	\vspace*{0.5cm}\noindent\textbf{Keywords}: Intellectual disability; Inclusive education; Artificial intelligence.
		
\end{ABSTRACT}
